\documentclass[pre,twocolumn,superscriptaddress,nofootinbib]{revtex4-2}
\usepackage{graphicx}
\usepackage{amsmath}
\usepackage{booktabs}
\usepackage[hidelinks]{hyperref}

\begin{document}

\title{What the burp leaves behind: a relic defect of fixed size in a graph model of emergent geometry, its
number, its position, and how long it lasts}

\author{Emily Smith}
\email{emilysmithcreate@gmail.com}
\affiliation{Independent researcher, Charlottesville, Virginia, USA}

\date{DRAFT, 25 September 2026}

\begin{abstract}
In a companion paper we showed that, in the family of graph energies obtained from combinatorial quantum
gravity (CQG) by raising the coefficient $\lambda$ of its local term above 1 (only $\lambda = 1$ is CQG; every
result here is at $\lambda > 1$, next to CQG and not in it), a $4 \times L$ torus with one direction curled is
metastable and opens into the flat torus by a single front, releasing exactly $\varepsilon = 4(\lambda - 1)$ per
vertex. Here we study what the change leaves behind. Sealed, with a bath cold enough to hold the release, the
conversion ends as a flat sheet with about one relic defect, at every size from 64 to 288 vertices: the relic does
not grow with the space. Read from its wiring it is a closed loop of four vertices at local dimension 1, costing
$24\lambda - 16$ above the flat torus in the cold boxes (14 at $\lambda = 1.25$), with lighter forms of 4 and 9 units
and rarer resting states of 8 to 45 units. Its position is neither at the seed nor opposite it: it is
indistinguishable from uniform. Planting $k = 1, 2, 4, 8$ seeds leaves $1.10, 2.15, 2.50, 3.40$ relics on
average: more seeds make more scrap, but fewer than one each. Held at a fixed coupling of 1.0, 1.25 or 1.5 the relic
anneals away within $500$ to $80{,}500$ sweeps; cooled from 1.25 toward a coupling at which nothing moves, it
freezes in: every relic survives a fast cooling, about four in five survive a cooling of ten to thirty thousand sweeps,
and a third survive the slowest tried. Every verdict quoted was pre-registered before its run.
\end{abstract}

\maketitle

\section{Introduction}

A change from one ordered arrangement to another releases energy and, in a sealed system, cannot end in the perfect
product: whatever the release cannot carry away stays behind as disorder. In the companion paper
\cite{paper1} (the curled torus) the change is a decompactification in miniature: a $4 \times L$ torus, one
direction curled to a single square, opens into the flat torus once the local term of CQG \cite{T17,KTB19,T25} is
weighted by a coefficient $\lambda > 1$. That paper measured the wait, the barrier, the front and the release. This
one measures the remainder: its size, its number, where it sits, and how long it lasts.

The question is narrow and the answer is measurable. The reason for asking is the author's hypothesis that a
change of this kind leaves a remainder, whatever did not convert, which has to account for the rest of the energy
books (claim 5 of the project's VISION). Ref.~\cite{T24} describes something in the same family for CQG itself: allotropes, regions of the network stuck in a different discrete arrangement, as diamond
survives where graphite is stable. The relic studied here is not an allotrope in that sense: it is a single object
of fixed size, made by the burp, at $\lambda > 1$. The two are cousins, and we say so where it matters.

Every kinetic statement below concerns one Markov chain, the local edge switch used to sample the model
\cite{KTB19}, with Metropolis acceptance, and the sealed version of it built from Creutz demons \cite{Creutz83}. Energies,
the relic's cost and the exact classification of end states do not depend on the chain; waiting times, annealing
times and the count of relics per seed do.

\section{Model, sealed dynamics, and how a relic is read}

The state is a four-regular bipartite graph on $N$ vertices obeying the hard-core rule of CQG (no two vertices share
more than two neighbors). With $S$ the number of squares and $X = \sum_e (S_e - 2)_+$ the number of surplus squares
on edges carrying more than two, the energy is
\begin{equation}
  H = 16(N - S) + 4\lambda X,
  \label{eq:H}
\end{equation}
which for $\lambda = 1$ is the total Ollivier curvature of CQG \cite{T25}, and for $\lambda \ne 1$ is not a curvature.
Written edge by edge, $H = 4\sum_e [(2 - S_e) + \lambda (S_e - 2)_+]$: the flat torus, with two squares on every
edge, sits at $H = 0$ and is the ground state for $\lambda \ge 1$; each curled direction costs $4(\lambda - 1)$ per
vertex above it.

The local dimension $d(v)$ of a vertex is the number of its six edge pairs that close no square: 2 on the flat sheet,
1 on the curled torus, 0 in a 4-cube. A relic is counted as a connected piece of vertices at $d = 1$ in an otherwise
flat sheet, and its energy above the flat torus is read exactly from Eq.~\eqref{eq:H} on the saved wiring. A defect
that leaves every vertex at $d = 2$ is not seen by this count; one such state, holding 20 units, is among the rare
states of Sec.~IV.

Sealed runs keep a bath of $C$ stores beside the graph. A proposed switch picks one store at random, is refused if
that store cannot pay for it, and otherwise pays or is paid; $H + \sum \text{stores}$ is conserved to the last unit
in every run reported here. The stores' mean energy reads the temperature the system reaches by itself; $C$ sets
how far that temperature rises per unit released, and the companion paper showed that at $\lambda = 1.25$ the change
ends as a clean sheet in most runs only when $C \ge N/2$ (the prediction, about $N/3.5$, lies between the grid points
$N/4$ and $N/2$). All runs below use $C = 2N$ stores, cold enough that the release is held and the sheet does not
melt: in T10 and T11 one store starts with a 12-unit spark and the rest are empty; in T17, T19 and T25 every store
starts empty and the change is started by planted seeds.

\section{One relic, however large the space}

\emph{Pre-registered as T10.} Twenty sealed conversions at each of $N = 64, 96, 192, 288$ ($4 \times L$ tori,
$\lambda = 1.25$, $C = 2N$, a 12-unit spark to start the change, 30,000 sweeps). The prediction on record was that
the mean number of relics rises with $N$; the verdict was ONE RING, HOWEVER LARGE: at every size the product is a
flat sheet with about one relic of four vertices, the mean count $0.85$ to $1.05$, within the replica scatter of 1,
with a slope of $0.0003 \pm 0.0003$ per vertex. A second check, that every relic holds the expected 14 units, failed
in 9 of the 80 boxes, which held instead a different leftover: a twist two points wide, costing 20 units. (The
pre-registered label keeps the word ``ring''; in these cold boxes the four points lie along the torus, spanning one
to four of its original columns, mostly two or three, not round its short direction.)

\section{What the relic is, read from its wiring}

Read from the saved final graphs of two of the companion paper's decays, the four-point relic is a closed loop of
four vertices at $d = 1$, each of its edges lying in three squares; in one of the two it is an original column of the
torus that stayed curled while the columns around it opened. At $\lambda = 1.25$ it comes in three forms, holding 4,
9 or 14 units according to how many of the edges around it carry a third square ($S = N + 1$ and $X = 4$, 5 or 6).
The form left in the cold boxes of T10 and T11 is the 14-unit one, costing $24\lambda - 16$ above the flat torus;
there it was identified from its energy and the positions of its points, not read from its wiring. Rarer resting
states, 8 to 45 units, are combinations of small defects in the sheet: a two-point twist (20 units), a 3-cube, an
open line of three, single edges and single points; eleven of twelve unusual states read this way matched their
recorded release exactly, and the twelfth was set aside as a decay that changed state inside its last measuring
window. In the states read these are small, 8 to 45 units at every tube length; whether the list of kinds is
complete is not known.

The share of the release kept by the 14-unit relic is exact: $(24\lambda - 16)/[4(\lambda - 1)N]$, which is 22\% at $\lambda = 1.25$
and $N = 64$, 72\% at $\lambda = 1.05$, and falls as $1/N$. Near $\lambda = 1$ the scrap keeps a large share of
the lump; in a large space it is negligible.

\section{Where it sits}

\emph{Pre-registered as T11}, in cold boxes as T10 at $N = 64$ and 96. Two readings were possible: the relic marks
the seed, or it marks the seam opposite the seed, where a front spreading both ways from one start would meet itself.
In the 36 and 28 boxes holding one relic its position is indistinguishable from uniform along the torus: neither.
Where fronts actually met was not measured. The verdict on record is NEITHER.

\section{How many: one per tube, or one per seed?}

\emph{Pre-registered as T17.} A $96 \times 4$ torus ($N = 384$) at $\lambda = 1.25$ in a cold sealed box, $k$ seeds
planted as the cheapest move out of the torus at evenly spaced columns, twenty runs per $k$, all converting, energy
exact. The author's hypothesis corresponds to one relic per seed, though she did not pick among the pre-registered
options; ours was between, leaning toward one per tube whatever $k$. Neither one per seed nor one per tube held:

\begin{center}
\begin{tabular}{@{}lcccc@{}}
\toprule
seeds $k$ & 1 & 2 & 4 & 8 \\
relics per tube & 1.10 & 2.15 & 2.50 & 3.40 \\
\bottomrule
\end{tabular}
\end{center}

\noindent The slope is $0.29 \pm 0.04$ relics per extra seed; the verdict on record is BETWEEN, though the slope is
only about one standard error above the boundary with one per tube (0.25), and a straight line describes poorly means
that rise by 1.05, then 0.35, then 0.90. More seeds make more scrap, but fewer than one each, and the growth looks
less than proportional. The counts are also more regular than independent events would give: their variance over
their mean is 0.09 to 0.50, against about 1. Where fronts from neighboring
seeds meet, their relics may merge or anneal; that was not measured and the saved end states hold the positions.

\section{How long it lasts at a fixed temperature}

\emph{Pre-registered as T19.} Sheets with exactly one relic, made as in T17 at $k = 1$ but from a $24 \times 4$ torus
($N = 96$), were held at fixed coupling
$g = 1.0$, $1.25$ and $1.5$ for 100,000 sweeps, the relic's position read every 500. The options on record were
that it moves, that it anneals at the warmer couplings and stays at the colder, or that it stays at every $g$. It
anneals at every coupling: at $g = 1.0$ all nine followed relics were gone between 4,000 and 80,500 sweeps; at 1.25
all ten within 500 to 13,000; at 1.5 eight annealed and two moved first, all gone within 11,000. In T10's cold boxes,
whose bath ends near $g = 0.5$, the relic lasted the whole 30,000 sweeps. Our own prediction, that it would stay at
$g = 1.0$, was wrong. Said plainly: in this model the relic did not last at any of the three fixed couplings tried,
and it lasted where its surroundings stayed very cold.

\section{The race: it freezes in when the box cools}

\emph{Pre-registered as T25.} Reality cools as it expands, so the question that matters for a long-lived relic is a
race between healing and cooling. A sheet with one relic, made as in Sec.~VII ($N = 96$), was cooled from
$g = 1.25$, where every
relic anneals within 13,000 sweeps at fixed coupling, to $g = 0.25$, where nothing moves, by the same factor each block
of 500 sweeps over $t_{\rm cool} = 300$ to $100{,}000$ sweeps, and then held cold for 20,000 sweeps; twenty replicas per
cooling time, of which 17 to 20 ended the making step with exactly one relic and were followed. The pre-registered
verdict was FREEZES IN if some cooling time keeps a majority of relics and the fastest cooling keeps the most, ALWAYS
HEALS if no cooling time keeps a majority. The prediction recorded for the author was FREEZES IN, with no
freeze-out time; it was inferred by the assistant from her stated positions, and she confirmed it as her own only
after the reading below had been recorded. Ours was FREEZES IN with a freeze-out time between 1,000 and 10,000
sweeps.

\begin{center}
\small
\begin{tabular}{@{}lcccccc@{}}
\toprule
$t_{\rm cool}$ (sweeps) & 300 & 1,000 & 3,000 & 10,000 & 30,000 & 100,000 \\
relics surviving & 20/20 & 17/17 & 19/19 & 15/19 & 14/18 & 7/19 \\
\bottomrule
\end{tabular}
\end{center}

\noindent The verdict is FREEZES IN, with the freeze-out time, the longest cooling that keeps a majority, at 30,000
sweeps. Even a cooling of 30,000 sweeps, more than twice the longest annealing time seen at the starting coupling in
T19, keeps most relics, because the coupling leaves the range where healing is fast early in the schedule: the relics
that healed were last seen within the first 15,000 sweeps of cooling, while the coupling was still above about 1. Our
prediction of the freeze-out time was too short by a factor of three to thirty; the verdict is the one recorded for
the author, whose prediction named no freeze-out time. A relic that anneals at every fixed coupling tried is
therefore kept by cooling in this model, which bears on whether a remainder of the change can last. What the model
cannot say is how its sweep relates to a cosmological clock.

\section{Discussion}

Four facts about the relic are now on the record for this model. It is small and of fixed size, a closed loop of
four points left curled, and it does not grow with the space. Its number depends on how the change started, growing
with the number of seeds but by less than one per seed. It anneals away at each fixed coupling tried, from 1.0 to
1.5. And it freezes in when its surroundings cool, even slowly.

For the author's hypothesis, which needs the change to leave a remainder, the first two say what that remainder is
in this model: small, of fixed size, and more of it when the change starts in more places. The last two say when
it lasts: a relic survives when its surroundings cool, even slowly, on a clock the model cannot relate to cosmic
time. What part such a remainder plays at cosmic scale is beyond what this model can address.

What this cannot show: anything about the real universe; how the model's cooling clock relates to a cosmological
one; and whether a relic in three dimensions, or with the model's points treated as interchangeable, behaves the
same way. Under interchangeable points the flat sheet carries $2N$ symmetries ($4N$ when it is square) and a sheet
with one relic a few, so
annealing should be favored by a factor of order $N$; that expectation is recorded, not measured.

\begin{acknowledgments}
Code, analysis and text were drafted with Claude (Anthropic) in conversation with the author, who directed and
checked the work; no physicist has reviewed it. The repository holds every configuration, seed, result and
pre-registration: \url{https://github.com/EmilySmithCreate/RecreationalPhysics}.
\end{acknowledgments}

\begin{thebibliography}{9}
\bibitem{paper1} E. Smith, The curled torus burps: front-driven decompactification between ordered states in a
graph model of emergent geometry, draft (2026), \texttt{docs/papers/curled\_torus/}.
\bibitem{T17} C. A. Trugenberger, Combinatorial quantum gravity: geometry from random bits, JHEP 09, 045 (2017).
\bibitem{KTB19} C. Kelly, C. A. Trugenberger and F. Biancalana, Self-assembly of geometric space from random
graphs, Class. Quantum Grav. 36, 125012 (2019).
\bibitem{T25} C. A. Trugenberger, Networks as the fundamental constituents of the universe, J. Phys. Complex. 6,
042001 (2025); arXiv:2512.17676.
\bibitem{T24} C. A. Trugenberger, Dark matter and dark energy in combinatorial quantum gravity, arXiv:2409.09385
(2024).
\bibitem{Creutz83} M. Creutz, Microcanonical Monte Carlo simulation, Phys. Rev. Lett. 50, 1411 (1983).
\end{thebibliography}

\end{document}
